\documentclass[10pt,letterpaper]{article}
\usepackage{cornell-notes}

\title{Chapter 2: Bipolar Transistors}
\author{}
\date{}
\setDocTitle{Chapter 2: Bipolar Transistors}
\setDocSubtitle{Electronics}
\setDocOwner{Electronics Cornell Notes}
\setDocDate{\today}
\setCornellCollection{Electronics Cornell Notes}
\setCornellUnitType{Chapter}
\setCornellUnitNumber{2}
\setCornellUnitTitle{Bipolar Transistors}
\hypersetup{pdftitle={Chapter 2: Bipolar Transistors},pdfsubject={Cornell notes},pdfauthor={}}
\begin{document}
\maketitle

\begin{CornellOverviewBox}{Concept}
Bipolar junction transistors (BJTs) are developed as switches, followers, current sources, voltage amplifiers, differential stages, and output drivers. Simple current-gain intuition is useful for switching, but predictable analog design comes from emitter degeneration, transconductance, feedback, and bias networks that reduce dependence on poorly controlled $\beta$ and temperature-sensitive $V_{BE}$.
\end{CornellOverviewBox}
\begin{CornellOverviewBox}{Concept}
\textbf{Prerequisites.} Chapter 1 dc networks, small-signal resistance, capacitors, power dissipation, and basic frequency response.

\textbf{After studying this chapter, you should be able to:}
\begin{itemize}
  \item Bias BJTs for switching and linear operation without relying on a precise $\beta$.
  \item Analyze emitter followers, common-emitter stages, current mirrors, and differential pairs.
  \item Use $g_m$, emitter degeneration, and negative feedback to predict gain and impedances.
  \item Evaluate output-stage crossover, sharing, stability, capacitance, and thermal limits.
\end{itemize}\textbf{Central design questions.}
\begin{itemize}
  \item Is the device intended to be cut off, saturated, or in forward-active operation?
  \item Which parameter spread or temperature effect dominates the bias point?
  \item How do feedback and degeneration trade gain for predictability, bandwidth, and linearity?
\end{itemize}\end{CornellOverviewBox}
\section{Cornell Notes}
\begin{CornellNotesTable}
\CornellNoteRow{2.1: First BJT model}{In forward-active operation, collector current is controlled exponentially by $V_{BE}$ and is approximately $\beta I_B$ only over a limited operating region. $\beta$ varies widely with device, current, voltage, and temperature. Use a base-current model to guarantee switch drive; use the transconductance model for analog increments. A silicon base-emitter drop near $0.6$--$0.8\,\mathrm{V}$ is a rough consequence, not a design constant.}
\CornellNoteRow{2.2: Basic circuits}{A transistor switch requires enough base drive for the worst-case load and driver while limiting junction stress. An emitter follower gives voltage gain near unity, current gain, low output impedance, and a level shift. A common-emitter stage converts input voltage to collector current and load voltage; emitter resistance stabilizes current and linearizes gain. Current sources need compliance voltage. Phase splitters generate opposite-polarity outputs.}
\CornellNoteRow{2.3: Ebers--Moll and $g_m$}{For small changes about a bias point, $g_m=I_C/V_T$ with $V_T\approx25.8\,\mathrm{mV}$ near room temperature; $r_e\approx V_T/I_E$. An unbypassed emitter resistor makes gain depend mainly on resistor ratio rather than transistor parameters. Current mirrors copy a reference current but suffer finite output resistance, base-current error, mismatch, and thermal gradients. Differential pairs steer a tail current and reject ideal common-mode input.}
\CornellNoteRow{2.4: Amplifier building blocks}{Push-pull stages source and sink load current efficiently but need bias control to limit crossover distortion and thermal runaway. Darlingtons provide high current gain at the cost of two junction drops, extra charge storage, and slower turn-off. Bootstrapping raises effective impedance by moving both ends of an element together. Paralleled BJTs require emitter ballast or active sharing because hotter devices tend to draw more current. Miller multiplication of base-collector capacitance limits bandwidth.}
\CornellNoteRow{2.5: Negative feedback}{With forward gain $A$ and feedback factor $f$, closed-loop gain is $A_f=A/(1+Af)$ for negative feedback under the assumed sign convention. Large loop gain reduces gain sensitivity, distortion, and some noise while changing input/output impedances and bandwidth. Stability is not automatic: phase lag can turn feedback positive. Preserve adequate phase margin and evaluate every high-impedance node and capacitive load.}
\CornellNoteRow{2.6: Complete circuits}{Regulators, thermal controllers, and transistor logic combine sensing, references, gain, and power stages. Partition the circuit into blocks, find the dc operating state, then trace small changes around the loop. Separate load-current paths from sensitive reference paths, provide current limiting, and check startup, shutdown, faults, and safe operating area rather than only nominal regulation.}
\end{CornellNotesTable}
\begin{CornellSummaryBox}{Summary}
BJTs are best treated as transconductance devices whose useful behavior is shaped by resistors and feedback. Switching design guarantees base drive and stress margins; linear design establishes a stable quiescent current, uses degeneration or feedback to control gain, and respects capacitance, output compliance, power, and thermal coupling. Parameter spread is expected and must be designed out rather than calibrated away casually.
\end{CornellSummaryBox}
\section{Worked Examples}
\begin{CornellExamBox}{Guaranteed relay switch}
A $24\,\mathrm{V}$ relay draws $80\,\mathrm{mA}$ and is switched by a $5\,\mathrm{V}$ logic output. Choose forced $\beta=10$, so $I_B=8\,\mathrm{mA}$. With $V_{BE}=0.8\,\mathrm{V}$, $R_B=(5-0.8)/8\,\mathrm{mA}=525\,\Omega$; choose $510\,\Omega$ after confirming the logic pin can source the current. At $V_{CE(sat)}=0.2\,\mathrm{V}$, transistor dissipation is about $16\,\mathrm{mW}$. Add a coil suppression path and verify turn-off time.
\end{CornellExamBox}
\begin{CornellExamBox}{Degenerated voltage gain}
A common-emitter stage carries $I_C=1\,\mathrm{mA}$, uses unbypassed $R_E=470\,\Omega$, and sees $R_C\parallel R_L=3.0\,\mathrm{k\Omega}$. Since $r_e\approx25.8\,\Omega$, the approximate gain is $A_v\approx-3000/(470+25.8)=-6.05$. The resistor dominates $r_e$, so gain is much less sensitive to current and temperature than an undegenerated stage.
\end{CornellExamBox}
\begin{CornellExamBox}{Feedback sensitivity}
An amplifier has open-loop gain $A=180$ and feedback factor $f=0.09$. Then $A_f=180/(1+16.2)=10.47$. A $20\%$ open-loop gain increase produces $A_f=10.56$, less than a $1\%$ change. This illustrates desensitization, but stability still requires checking phase where $|Af|$ approaches unity.
\end{CornellExamBox}
\section{Design and Troubleshooting}
\begin{CornellExamBox}{Design and Troubleshooting}
\begin{enumerate}
  \item Choose operating region, load current/voltage, supply, required gain, bandwidth, and allowed error.
  \item Set a bias point with compliance and power margin; use degeneration or feedback to reduce dependence on $\beta$ and $V_{BE}$.
  \item Evaluate small-signal gain, input/output resistance, capacitances, noise, and stability.
  \item Check startup, cutoff, saturation recovery, short circuits, temperature, safe operating area, and component tolerances on hardware.
\end{enumerate}\end{CornellExamBox}
\begin{CornellExamBox}{Symptoms, likely causes, and checks}
\begin{itemize}
  \item Collector stuck high: missing base drive, open emitter path, wrong pinout, or insufficient reference return.
  \item Collector stuck low: excessive base drive, unintended saturation, leakage, reversed device, or load fault.
  \item Bias drifts with heat: poor emitter degeneration, shared thermal path, inadequate heatsinking, or runaway.
  \item Low high-frequency gain: Miller capacitance, source resistance, wiring capacitance, or insufficient bias current.
  \item Oscillation with a load: loop phase shift, capacitive loading, poor decoupling, or long feedback routing.
\end{itemize}\end{CornellExamBox}
\begin{CornellWarningBox}{Safety and Limitations}
Power BJTs can fail from secondary breakdown even when average power seems acceptable; use the datasheet safe-operating-area graph at the actual pulse duration and temperature. Inductive loads require controlled energy release. Use current-limited supplies during bring-up, and assume heatsinks and exposed collectors may be electrically live.
\end{CornellWarningBox}
\section{Key-Equation Sheet}
\begin{CornellExamBox}{Key Equations}
\begin{itemize}
  \item $I_C\approx\beta I_B$ is a rough forward-active relation; switching uses a conservative forced $\beta$.
  \item $g_m=I_C/V_T$ and $r_e\approx V_T/I_E$, with $V_T\approx25.8\,\mathrm{mV}$ near room temperature.
  \item Degenerated common-emitter gain: $A_v\approx-(R_C\parallel R_L)/(R_E+r_e)$ when other loading is negligible.
  \item Follower output resistance is roughly the base-source resistance divided by $\beta+1$, in parallel with emitter-network resistance.
  \item $A_f=A/(1+Af)$ for a negative-feedback loop under the stated sign convention.
  \item $P_Q=V_{CE}I_C$ at the quiescent point; junction temperature needs thermal-resistance analysis.
\end{itemize}\end{CornellExamBox}
\section{Study and Review}
\subsection{Design checklist}
\begin{itemize}
  \item Worst-case $\beta$, driver current, and saturation recovery are acceptable.
  \item Bias current and $V_{CE}$ provide signal swing and compliance.
  \item Junction temperature and safe operating area pass normal and fault cases.
  \item Feedback polarity, loop gain, decoupling, and capacitive loading are checked.
  \item Pinout, package tab connection, and complementary-device differences are verified.
\end{itemize}\subsection{Common misconceptions}
\begin{itemize}
  \item $V_{BE}$ is not exactly $0.7\,\mathrm{V}$.
  \item A high catalog $\beta$ does not guarantee saturation with tiny base current.
  \item Saturation is useful for a switch but harmful to fast linear recovery.
  \item Parallel BJTs do not automatically share current.
  \item Negative feedback does not guarantee stability.
\end{itemize}\subsection{Active-recall questions}
\begin{enumerate}
  \item Name the three common BJT operating regions.
  \item Why use forced $\beta$ for a switch?
  \item What does emitter degeneration improve?
  \item How does $g_m$ scale with collector current?
  \item What limits current-source compliance?
  \item Why does a differential pair reject common-mode input ideally?
  \item What causes crossover distortion?
  \item Why can paralleling cause thermal imbalance?
  \item What is the Miller effect?
  \item Which condition determines feedback stability?
\end{enumerate}\subsection{Original practice problems}
\begin{enumerate}
  \item A $3.3\,\mathrm{V}$ driver must switch $40\,\mathrm{mA}$. With forced $\beta=8$ and $V_{BE}=0.8\,\mathrm{V}$, find $R_B$.
  \item Estimate $g_m$ and $r_e$ at $I_C\approx I_E=2\,\mathrm{mA}$.
  \item A device dissipates $0.8\,\mathrm{W}$ with total junction-to-ambient thermal resistance $70^{\circ}\mathrm{C/W}$ at $30^{\circ}\mathrm{C}$ ambient. Estimate junction temperature.
\end{enumerate}\subsection{Answer key / solution outline}
\begin{enumerate}
  \item $I_B=5\,\mathrm{mA}$ and $R_B=(3.3-0.8)/5\,\mathrm{mA}=500\,\Omega$; verify driver capability and choose a standard value conservatively.
  \item $g_m=0.002/0.0258=77.5\,\mathrm{mS}$ and $r_e\approx12.9\,\Omega$.
  \item Rise is $56^{\circ}\mathrm{C}$, so $T_J\approx86^{\circ}\mathrm{C}$ before considering transient or local heating.
\end{enumerate}\begin{CornellSummaryBox}{Summary}
Use $\beta$ only as a bounded convenience. Set current and gain with resistors, degeneration, and feedback; guarantee switching drive; and verify compliance, capacitance, thermal behavior, and safe operating area under real load and fault conditions.
\end{CornellSummaryBox}

\end{document}
